\documentclass[border=14pt]{standalone}
\usepackage{ctex}
\usepackage{tikz}
\begin{document}
\begin{tikzpicture}[x=1cm, y=1cm, font=\normalsize, line width=0.8pt]

% fig-ge-02 相量加减 = 首尾相接（中文；页面算例 U1=3+j4、U2=3-j4）
% 结构约定（前三版都栽在线穿字上，这版按位置分区）：
%   左栏：纯文字，无任何横线分隔（分隔线会压字，删掉）
%   右栏：每半幅 = 左侧坐标系（轴止于 x=4.0）+ 右侧文字块（x>=4.6），
%         两区之间留 0.6 空隙，任何文字都不进坐标系区域

\def\colw{6.6}
\def\gap{2.0}

% ── 左栏：代数形式 ──
\begin{scope}[shift={(0,0)}]
  \node[anchor=south] at ({\colw/2},6.4) {代数形式};
  \node[anchor=north west, font=\small] at (0.4,6.0) {U1 = 3 + j4};
  \node[anchor=north west, font=\small] at (0.4,5.45) {U2 = 3 - j4};
  \node[anchor=north west, font=\small, align=left] at (0.4,4.6)
    {和\\ 实部 3+3 = 6\\ 虚部 4+(-4) = 0\\ 和 = 6 + j0};
  \node[anchor=north west, font=\small, align=left] at (0.4,2.6)
    {差\\ 实部 3-3 = 0\\ 虚部 4-(-4) = 8\\ 差 = 0 + j8};
  \node[anchor=north west, font=\small] at (0.4,0.7)
    {虚部一正一负：和抵消、差加倍};
\end{scope}

% ── 右栏：相量图（上和 / 下差）──
\begin{scope}[shift={({\colw+\gap},0)}]
  \node[anchor=south] at ({\colw/2},6.4) {相量图};

  % 上半：和
  \begin{scope}[shift={(0,3.5)}]
    \draw[->] (1.0,1.0) -- (4.0,1.0);
    \draw[->] (1.0,0.1) -- (1.0,3.0);
    \draw[->, line width=1pt] (1.0,1.0) -- ++(1.05,1.4);      % U1
    \draw[->, line width=1pt] (1.0,1.0) -- ++(1.05,-1.4);     % U2
    \draw[dashed, gray, line width=0.5pt] (2.05,2.4) -- ++(1.05,-1.4);
    \draw[dashed, gray, line width=0.5pt] (2.05,-0.4) -- ++(1.05,1.4);
    \draw[->, line width=1.7pt] (1.0,1.0) -- ++(2.1,0);       % 和
    \node[anchor=east, font=\small] at (0.90,1.0) {Im};
    \node[anchor=west, font=\small] at (4.05,1.0) {Re};
    % 文字块：x>=4.6，与坐标系不重叠
    \node[anchor=north west, font=\small, align=left] at (4.65,2.9) {U1\\ U2};
    \node[anchor=north west, font=\small] at (4.65,1.5) {和 = 6 + j0};
  \end{scope}

  % 下半：差
  \begin{scope}[shift={(0,-0.2)}]
    \draw[->] (1.0,1.0) -- (4.0,1.0);
    \draw[->] (1.0,0.1) -- (1.0,3.0);
    % U1 指向右上 (+3,+4)；-U2 = -3+j4 指向左上 (-3,+4)，首尾相接后合力竖直向上
    \draw[->, dashed, line width=0.9pt] (1.0,1.0) -- ++(0.50,0.66);      % U1 → 右上
    \draw[->, dashed, line width=0.9pt] (1.5,1.66) -- ++(-0.50,0.66);    % -U2 → 左上（从 U1 末端起）
    \draw[->, line width=1.7pt] (1.0,1.0) -- ++(0,1.32);                 % 差 = 0+j8 竖直
    \node[anchor=west, font=\small] at (1.12,2.05) {差};
    \node[anchor=north east, font=\small] at (0.30,0.62) {U1};
    \node[anchor=south east, font=\small] at (0.42,1.70) {-U2};
    \node[anchor=east, font=\small] at (0.90,1.0) {Im};
    \node[anchor=west, font=\small] at (4.05,1.0) {Re};
    \node[anchor=north west, font=\small, align=left] at (4.65,2.9) {U1\\ -U2（首尾相接）};
    \node[anchor=north west, font=\small] at (4.65,1.5) {差 = 0 + j8};
  \end{scope}
\end{scope}

\end{tikzpicture}
\end{document}